\begin{exercise}{1.1}{4}
    Man zeige: Die Gruppe $D \subset \operatorname{GL}_{\mathbb R}(\mathbb C)$ aller $\mathbb R$-linearen Automorphismen von $\mathbb C$ der Gestalt $w \mapsto z \cdot w$ oder $w \mapsto z \cdot \overline{w}$ für $z \in \mathbb C$ mit $|z| = 1$
    ist eine Drehspiegelgruppe. Man gebe dafür ein invariantes Skalarprodukt auf
    dem $\mathbb R$-Vektorraum $\mathbb C$ an.
\end{exercise}
\noindent
To show $D$ is contained in $O(2)$, we need to check that it consists
of isometries with respect to some inner product. For the $\mathbb R$-vector space $\mathbb C$,
let us define the inner product as: $\langle u, v \rangle = \operatorname{Re}(u\overline{v})$.
Note that if we let $u = u_1 + iu_2$ and $v = v_1 + iv_2$, this yields the usual dot
product $u_1v_1 + u_2v_2$.
Let $u, v \in \mathbb C$. For automorphisms of the first type, we check
\[
    \langle zu, zv \rangle = \operatorname{Re}(zu \cdot z\overline{v}) = \operatorname{Re}(z\overline{z} \cdot u\overline{v}) = \operatorname{Re}(u\overline{v}) = \langle u, v\rangle
\]
For automorphisms of the second type, we check
\[
    \langle z\overline{u}, z\overline{v} \rangle = \operatorname{Re}(z\overline{u} \cdot \overline{z\overline{v}}) = \operatorname{Re}(z\overline{z} \cdot u\overline{v}) = \operatorname{Re}(\overline{u}v) = \operatorname{Re}(u\overline{v}) = \langle u, v\rangle
\]
So indeed $D \subseteq O(2)$.
The subgroup $SO(2)$ of $O(2)$ consists of all rotations. As we let $z$ run over
all elements of $\mathbb C$ with modulus 1, writing $z = e^{i\theta}$ we see that we obtain all
rotations by an angle $\theta \in [0, 2\pi]$. The full group $O(2)$ is generated by $SO(2)$
along with the reflection $w \mapsto \overline{w}$, so indeed $D = O(2)$.

\textit{Zusammenfassung:}
\begin{enumerate}
    \item Zeige zuerst $D \subseteq O(2)$ mit gewähltem Skalarprodukt $\langle u, v \rangle = \operatorname{Re}(u\overline{v})$
    \item $\langle zu, zv \rangle = \operatorname{Re}(zu \cdot z\overline{v}) = \operatorname{Re}(z\overline{z} \cdot u\overline{v}) = \operatorname{Re}(u\overline{v}) = \langle u, v\rangle$ weil $z\overline{z} = 1$, da $|z| = 1$
    \item $\langle z\overline{u}, z\overline{v} \rangle = \operatorname{Re}(z\overline{u} \cdot \overline{z\overline{v}}) = \operatorname{Re}(z\overline{z} \cdot u\overline{v}) = \operatorname{Re}(\overline{u}v) = \operatorname{Re}(u\overline{v}) = \langle u, v\rangle$
    \item Nun zeige, dass $D = O(2)$, indem man $SO(2)$ und $O(2)\setminus SO(2)$ erzeugt
    \item $SO(2)$ wird erzeugt mit $z = e^{i\theta} \operatorname{mod} 1$ und $\theta = [0, 2\pi]$
    \item $O(2) \setminus SO(2)$ mit $w \mapsto \overline{w}$
\end{enumerate}

\begin{exercise}{1.2}{3}
    Man zeige: Gegeben $(Z, D)$ ein zweidimensionaler reeller Vektorraum mit ausgezeichneter Drehspiegelgruppe wird die Drehspiegelgruppe $D$
    von ihren Spiegelungen erzeugt.
\end{exercise}
Sei $S \subset D$ die Menge aller Spiegelungen in D. Z.z.: $D = \langle S\rangle$.
Da $S \subset D$ und $D$ eine Gruppe ist, gilt $\langle S\rangle \subset D$.
Sei $r \in D$ eine Drehung und $s \in D$ eine Spiegelung. Dann ist $rs$ wieder eine Spiegelung, weil
$\operatorname{det}(rs) = \operatorname{det}(r) \cdot \operatorname{det}(s) = 1 \cdot (-1) = -1$.
Da für jede Spiegelung gilt $s^2 = \operatorname{id}$ folgt $r = r\cdot \operatorname{id} = rs^2 = (rs)s \in \langle S\rangle$.
Da $rs \in S$ und $s \in S$, folgt $r \in \langle S \rangle$, also $D \subset \langle S\rangle$ und somit $D = \langle S\rangle$.

\textit{Zusammenfassung:}
\begin{enumerate}
    \item $S \subset D$ Menge aller Spiegelungen. Z.Z.: $D = \langle S\rangle$
    \item $\langle S\rangle \subseteq D$, weil $S \subset D$ und $D$ ist Gruppe
    \item Nun $r \in D$ Drehung, $s \in D$ Spiegelung, dann $rs \in S$ ($\operatorname{det}(rs) = \operatorname{det}(r) \cdot \operatorname{det}(s)$)
    \item Es gilt $s^2 = \operatorname{id}$ für alle Spiegelungen, also $r = r\cdot \operatorname{id} = rs^2 = (rs)s$
    \item Weil $rs, s \in S$ folgt $r \in \langle S\rangle$ und damit $D \subseteq \langle S \rangle$.
\end{enumerate}

%\begin{exercise}{1.3}{1}
%    Man bestimme alle endlichen Untergruppen von $O(2)$.
%\end{exercise}

\begin{exercise}{1.4}{3}
    Seien \((Z,s,\varepsilon)\) und \((U,t,\eta)\) zweidimensionale reelle
    Vektorräume mit Skalarprodukt und Orientierung. Für einen orthogonalen
    Isomorphismus $\psi:Z\xrightarrow{\sim} U$
    erklären wir
    \[
    \operatorname{int}_{\psi}:SO(Z,s)\xrightarrow{\sim}SO(U,t)
    \]
    durch $k\mapsto \psi\circ k\circ \psi^{-1}.$
    Man zeige, dass für je zwei orthogonale orientierungserhaltende Abbildungen
    $\varphi,\psi:Z\xrightarrow{\sim} U$
    gilt:
    $\operatorname{int}_{\psi}=\operatorname{int}_{\varphi}.$
\end{exercise}

Es ist zu zeigen, dass für alle \(k\in SO(Z,s)\) gilt
$
\psi\circ k\circ \psi^{-1}
=
\varphi\circ k\circ \varphi^{-1}.
$
Dies ist äquivalent zu
$
\varphi^{-1}\circ \psi\circ k\circ \psi^{-1}\circ \varphi
=
k.
$
Setze $g:=\varphi^{-1}\circ \psi.$
Da \(\varphi\) und \(\psi\) orthogonal und orientierungserhaltend sind, ist auch
\(g\) orthogonal und orientierungserhaltend, also gilt $g\in SO(Z,s)$.
Au\ss erdem ist
$g^{-1}=\psi^{-1}\circ \varphi.$
Damit genügt es zu zeigen, dass
$g\circ k\circ g^{-1}=k$
für alle \(k\in SO(Z,s)\).
Da \(Z\) zweidimensional ist, ist \(SO(Z,s)\) kommutativ. Also gilt
$g\circ k=k\circ g.$
Daher folgt
$
g\circ k\circ g^{-1}
=
k\circ g\circ g^{-1}
=
k.
$

\textit{Zusammenfassung:}
\begin{enumerate}
    \item z.Z.: Für alle \(k\in SO(Z,s)\) gilt $\psi\circ k\circ \psi^{-1}=\varphi\circ k\circ \varphi^{-1}$
    \item Das ist äquivalent zu $\varphi^{-1}\circ \psi\circ k\circ \psi^{-1}\circ \varphi=k$
    \item Setze $g :=\varphi^{-1}\circ \psi$. 
    \item \(\varphi\) und \(\psi\) orthogonal und orientierungserhaltend, also auch $g$ und $g \in SO(Z,s)$
    \item Bemerke, dass $g^{-1}=\psi^{-1}\circ \varphi$
    \item Bemerke, dass $SO(Z,s)$ kommutativ, also $g\circ k=k\circ g$
    \item Es folgt $g\circ k\circ g^{-1}=k\circ g\circ g^{-1}=k$
\end{enumerate}